%%%%%%%%%%%%%%%%%%%%%%%%%%%%%%%%%%%%%%%%%%%%%%%%%%%%%%%%%%%%%%%%%%%%%%
% Lớp 10 | Thống kê | Các số đặc trưng đo xu thế trung tâm của mẫu số liệu
% Dạng 0D6?3-1: Câu hỏi lý thuyết
% Số câu: 14  (sinh tự động bởi tach_cau.py; từ nay thầy sửa trực tiếp file này)
%%%%%%%%%%%%%%%%%%%%%%%%%%%%%%%%%%%%%%%%%%%%%%%%%%%%%%%%%%%%%%%%%%%%%%

%%%%%%%%%%%%%%%% TRẮC NGHIỆM 4 PHƯƠNG ÁN %%%%%%%%%%%%%%%%

% nguồn: Ok_TracNghiemLop11_Lan1
\begin{ex}[c]%[0D6B3-1]%[ma:f9a4261a]
	Giá trị của biểu thức $\cos  \dfrac{37\pi}{12}$ bằng
	 \choice
		{\True  $-\dfrac{\sqrt{6}+\sqrt{2}}{4}$}
		 {$\dfrac{\sqrt{2}-\sqrt{6}}{4}$}
		{$\dfrac{\sqrt{6}+\sqrt{2}}{4}$} 
		 {$\dfrac{\sqrt{6}-\sqrt{2}}{4}$}
	\loigiai{ 
	Ta có $$\cos \dfrac{37\pi}{12}=\cos \left(\dfrac{3\pi}{4}+\dfrac{\pi}{3}  \right)=-\dfrac{\sqrt{6}+\sqrt{2}}{4}$$}
\end{ex}

% nguồn: Ok_TracNghiemLop11_Lan1
\begin{ex}%[0D6K3-1]%[ma:1461b3b8]
	Kết quả rút gọn của biểu thức $ \left( \dfrac{\sin{\alpha}+\tan{\alpha}}{\cos{\alpha}+1} \right)^2+1 $ bằng 
	 \choice
		{$ 2  $}{$ 1+\tan{\alpha} $}{\True $ \dfrac{1}{{{\cos}^2}{\alpha}} $}{$\dfrac{1}{{{\sin}^2}{\alpha}}$}
	\loigiai{Ta có  \\
	$ 
	\begin{aligned}
		\left(  \dfrac{\sin{\alpha}+\tan{\alpha}}{\cos{\alpha}+1} \right)^2+1 &= \left[\dfrac{\sin{\alpha}+\dfrac{\sin{\alpha}}{\cos{\alpha}}}{\cos{\alpha}+1}\right]^2+1 = \left[\dfrac{\sin{\alpha}(1+\cos{\alpha})}{\cos{\alpha}(\cos{\alpha}+1)}\right]^2+1 \\
		&= \tan^2\alpha +1 =  \dfrac{1}{\cos^2\alpha}.
	\end{aligned}
	 $
	 }
\end{ex}

% nguồn: Ok_TracNghiemLop11_Lan1
\begin{ex}[c]%[0D6K3-1]%[ma:1d7ea41a]
	Biết $\sin \alpha+\cos \alpha=m$. Tính $  P=\cos \left(\alpha-\dfrac{\pi}{4}\right)$ theo $m$
	 \choice
		{\True $ \dfrac{m}{\sqrt{2}}  $}
		 {$ m\sqrt{2} $ }
		 {$\dfrac{m}{2}$ }
		{$ 2m  $ }
	\loigiai{Ta  có $P=\cos \left(\alpha-\dfrac{\pi}{4}\right)=\cos \alpha \cdot \sin \dfrac{\pi}{4}+\sin \alpha \cos \dfrac{\pi}{4}=\dfrac{1}{\sqrt{2}} \cos \alpha+\dfrac{1}{\sqrt{2}} \sin \alpha$\\
		Do đó $ P=\dfrac{1}{\sqrt{2}}(\sin  \alpha+\cos \alpha)=\dfrac{m}{\sqrt{2}}$.
	} 
\end{ex}

% nguồn: Ok_TracNghiemLop11_Lan1
\begin{ex}%[0D6B3-1]%[ma:975bc016]
	Nếu biết $\sin \alpha =\dfrac{5}{13}\left(\dfrac{\pi }{2}<\alpha <\pi \right),\cos  \beta =\dfrac{3}{5}\left(0<\beta <\dfrac{\pi }{2}\right)$ thì giá trị đúng của $\cos (\alpha -\beta)$ là
	\choice 
		{ $\dfrac{16}{65}$} 
		{ \True $-\dfrac{16}{65}$} 
		{  $\dfrac{18}{65}$}
		{  $-\dfrac{18}{65}$}
	\loigiai{ 
		$\sin  \alpha =\dfrac{5}{13}\left(\dfrac{\pi }{2}<\alpha <\pi \right)\Rightarrow \cos \alpha =-\sqrt{1-\dfrac{25}{169}}=-\dfrac{12}{13}$.\\
		$\cos  \beta =\dfrac{3}{5}\left(0<\beta <\dfrac{\pi }{2}\right)\Rightarrow \sin \beta =\sqrt{1-\dfrac{9}{25}}=\dfrac{4}{5}$.\\
		Do đó $\cos (\alpha -\beta)=\cos \alpha \cdot \cos \beta  +\sin \alpha \cdot \sin \beta =-\dfrac{12}{13}\cdot \dfrac{3}{5}+\dfrac{5}{13}\cdot \dfrac{4}{5}=-\dfrac{16}{65}$.
	 }
\end{ex}

% nguồn: Ok_TracNghiemLop11_Lan1
\begin{ex}[c]%[0D6K3-1]%[ma:3d401cfb]
	Cho $\sin \alpha =\dfrac{1}{\sqrt{3}}$ với $0<\alpha<\dfrac{\pi}{2}$. Giá trị của $\cos \left(\alpha+\dfrac{\pi}{3}\right)$  bằng
	\choice 
		{$\dfrac{2-\sqrt{6}}{2\sqrt{6}}$} 
		 {$\sqrt{6}-\dfrac{1}{2}$}
		 { $\sqrt{6}-3$}
		 {\True $\dfrac{1}{\sqrt{6}}-\dfrac{1}{2}$}
	\loigiai{ 
		Ta có $\sin^2 \alpha+\cos^2 \alpha =1 \Rightarrow \cos^2  \alpha =1-\sin^2 \alpha =1-\left(\dfrac{1}{\sqrt{3}}\right)^2 = \dfrac{2}{3}$. Suy ra $\cos \alpha =\pm \dfrac{\sqrt{6}}{3}$.\\
		Vì $0<\alpha<\dfrac{\pi}{2}$ nên điểm  cuối cung $\alpha$ thuộc góc phần tư thứ $I$ nên $\cos \alpha >0$. Vậy $\cos \alpha = \dfrac{\sqrt{6}}{3}$.\\
		Khi  đó\allowdisplaybreaks
		\begin{eqnarray*}
			\cos \left(\alpha+\dfrac{\pi}{3}\right)&= &\cos \alpha  \cos \dfrac{\pi}{3}-\sin \alpha\sin \dfrac{\pi}{3}\\
			&=& \dfrac{\sqrt{6}}{3}\cdot  \dfrac{1}{2}-\dfrac{1}{\sqrt{3}}\cdot \dfrac{\sqrt{3}}{2}\\
			&=& \dfrac{\sqrt{6}}{6}-\dfrac{1}{2}=\dfrac{1}{\sqrt{6}}-\dfrac{1}{2}. 
		\end{eqnarray*}
	} 
\end{ex}

% nguồn: Ok_TracNghiemLop11_Lan1
\begin{ex}%[0D6K3-1]%[ma:7f38bb19]
	Cho $  \cos \alpha = \dfrac{3}{5}$ $\left( - \dfrac{\pi}{2} < \alpha < 0 \right) $. Tính giá trị của $ \sin \left(\dfrac{\pi}{6} - \alpha \right) $.
	 \choice
		{\True $ \dfrac{3  - 4\sqrt{3}}{10} $}
		 {$ \dfrac{4 + 3\sqrt{3}}{10} $}
		{$ \dfrac{3 + 4\sqrt{3}}{10} $} 
		{$ \dfrac{4 -  3\sqrt{3}}{10} $}
	\loigiai{ 
		Do $ - \dfrac{\pi}{2} < \alpha < 0 $ $  \Rightarrow \sin \alpha < 0 $. \\
		Ta có  $ \sin \alpha = - \sqrt{1 - \cos^2\alpha} = - \sqrt{1 - \dfrac{9}{25}} = - \dfrac{4}{5} $.\\
	Vậy $ \sin \left(\dfrac{\pi}{6} - \alpha \right) = \sin \dfrac{\pi}{6}\cos \alpha - \sin \alpha\cos \dfrac{\pi}{6} = \dfrac{1}{2}\cdot \dfrac{3}{5} - \dfrac{4}{5}\cdot \dfrac{\sqrt{3}}{2} = \dfrac{3 - 4\sqrt{3}}{10}  $.}
\end{ex}

% nguồn: Ok_TracNghiemLop11_Lan1
\begin{ex}[c]%[0D6B3-1]%[ma:60cb9019]
	Biểu thức  $\sin \left(a+\dfrac{\pi}{6}\right)$ được viết lại
	 \choice
		{$\sin \left(a+\dfrac{\pi}{6}\right)=\sin  a+\dfrac{1}{2}$}
		{$\sin \left(a+\dfrac{\pi}{6}\right)=\dfrac{1}{2}\sin a-\dfrac{\sqrt{3}}{2}\cos a$} 
		{$\sin \left(a+\dfrac{\pi}{6}\right)=\dfrac{\sqrt{3}}{2}\sin a-\dfrac{1}{2}\cos  a$}
		 {\True $\sin \left(a+\dfrac{\pi}{6}\right)=\dfrac{\sqrt{3}}{2}\sin a+\dfrac{1}{2}\cos a$}
	 \loigiai{
	Ta có $\sin \left(a+\dfrac{\pi}{6}\right)=\sin a\cdot\cos \dfrac{\pi}{6}+\cos a\cdot \sin\dfrac{\pi}{6}=\dfrac{1}{2}\cos a +\dfrac{\sqrt{3}}{2}\sin a$.} 
\end{ex}

% nguồn: Ok_TracNghiemLop11_Lan1
\begin{ex}[c]%[0D6B3-1]%[ma:8fe7fe94]
	Giá trị của biểu thức $\dfrac{\cos \dfrac{\pi}{10} \cos \dfrac{\pi}{15}-\sin \dfrac{\pi}{15} \sin \dfrac{\pi}{10}}{\cos \dfrac{\pi}{5} \cos \dfrac{2 \pi}{15}-\sin \dfrac{2 \pi}{15} \sin  \dfrac{\pi}{5}}$ bằng
	 \choice
		{\True$ \sqrt{3} $ } 
		 { $ 1 $}
		{$ -1 $  }
		{ $  \dfrac{1}{2} $}
	 \loigiai{
	Ta có: $\dfrac{\cos \dfrac{\pi}{10} \cos \dfrac{\pi}{15}-\sin \dfrac{\pi}{15} \sin \dfrac{\pi}{10}}{\cos \dfrac{\pi}{5} \cos \dfrac{2 \pi}{15}-\sin \dfrac{2 \pi}{15} \sin \dfrac{\pi}{5}}=\dfrac{\cos \left(\dfrac{\pi}{10}+\dfrac{\pi}{15}\right)}{\cos \left(\dfrac{2 \pi}{15}+\dfrac{\pi}{5}\right)}=\dfrac{\cos \dfrac{\pi}{6}}{\cos \dfrac{\pi}{3}}=\dfrac{\sin \dfrac{\pi}{3}}{\cos \dfrac{\pi}{3}}=\tan  \dfrac{\pi}{3}=\sqrt{3}$.}
\end{ex}

% nguồn: Ok_TracNghiemLop11_Lan1
\begin{ex}[c]%[0D6B3-1]%[ma:25b16efb]
	Cho $x$, $y$ là các góc nhọn và dương  thỏa $\cot x=\dfrac{3}{4},\cot y=\dfrac{1}{7}$. Tổng $ x+y$ bằng
	\choice 
		{$\dfrac{\pi }{3}$} 
		 {$\pi $}
		 {$\dfrac{\pi }{4}$}
		{\True $\dfrac{3\pi  }{4}$}
	 \loigiai{
		Ta có: $\cot (x+y)=\dfrac{\cot x\cdot \cot y-1}{\cot x+\cot y}=-1$ $\Rightarrow $ $ x+y=\dfrac{3\pi }{4}$ (Do  $ x,y$ là các góc nhọn và dương).
	 }
\end{ex}

% nguồn: Ok_TracNghiemLop11_Lan1
\begin{ex}[c]%[0D6B3-1]%[ma:75a34613]
	Cho các góc lượng giác $a,b$  và $T=\cos (a+b)\cos (a-b)-\sin (a+b)\sin (a-b)$. Mệnh đề sau đây đúng?
	 \choice
		 {$T=\sin 2b$}
		 {\True $T=\cos 2a$}
		{$T=\sin  2a$}
		{$T=\cos 2b$} 
	 \loigiai{
	Ta có $T=\cos (a+b)\cos (a-b)-\sin (a+b)\sin (a-b)=\cos [(a+b)+(a-b)]=\cos 2a$.} 
\end{ex}

% nguồn: Ok_TracNghiemLop11_Lan1
\begin{ex}%[0D6K3-1]%[ma:744e4f9a]
	Cho $\sin \alpha =\dfrac{5}{13}$ và $0<\alpha  <\dfrac{\pi}{2}$, khi đó giá trị của $\cos \left(\alpha -\dfrac{\pi}{4}\right)$ là
	\choice 
		{$\dfrac{\sqrt{2}}{34}$} 
		{$-\dfrac{\sqrt{2}}{26}$} 
		{\True $\dfrac{17\sqrt{2}}{26}$} 
		{$-\dfrac{7\sqrt{2}}{26}$} 
	\loigiai{ 
		Do $0<\alpha <\dfrac{\pi}{2}$ nên $\cos\alpha >0$.  Bởi vậy:
		$$\cos\alpha  =\sqrt{1-\sin^2\alpha}=\sqrt{1-\left(\dfrac{5}{13}\right)^2}=\dfrac{12}{13}.$$
		 Do đó:\allowdisplaybreaks
		\begin{align*}
			 \cos \left(\alpha -\dfrac{\pi}{4}\right)&=\cos\alpha\cos\dfrac{\pi}{4}+\sin\alpha\sin\dfrac{\pi}{4}\\
			&=\dfrac{12}{13}\cdot  \dfrac{\sqrt{2}}{2}+\dfrac{5}{13}\cdot \dfrac{\sqrt{2}}{2}=\dfrac{17\sqrt{2}}{26}.
		\end{align*}
	 }
\end{ex}

% nguồn: Ok_TracNghiemLop11_Lan1
\begin{ex}%[0D6B3-1]%[ma:87df9340]
	Cho các góc $\alpha $, $\beta $ thỏa mãn $\dfrac{\pi}{2}<\alpha $, $\beta <\pi $, $\sin \alpha =\dfrac{1}{3}$,  $\cos \beta =-\dfrac{2}{3}$. Tính $\sin \left(\alpha +\beta \right)$.
	 \choice
		{\True $\sin \left(\alpha +\beta  \right)=-\dfrac{2+2\sqrt{10}}{9}$}
		{$\sin \left(\alpha +\beta \right)=\dfrac{2\sqrt{10}-2}{9}$} 
		{$\sin \left(\alpha +\beta \right)=\dfrac{\sqrt{5}-4\sqrt{2}}{9}$} 
		{$\sin  \left(\alpha +\beta \right)=\dfrac{\sqrt{5}+4\sqrt{2}}{9}$}
	\loigiai{ 
		 Do $\dfrac{\pi}{2}<\alpha $, $\beta <\pi $ $ \Rightarrow \left\{
		\begin{aligned}
			 &\cos \alpha <0\\
			&\sin \beta >0.\\ 
		\end{aligned}
		\right. $\\ 
		Ta có $\cos \alpha =-\ \sqrt{1-\sin^2\alpha}=-\ \sqrt{1-\dfrac{1}{9}}=-\ \dfrac{2\sqrt{2}}{3}$. $\sin \beta  =\sqrt{1-\cos^2\beta}=\sqrt{1-\dfrac{4}{9}}=\dfrac{\sqrt{5}}{3}$.\\
		Suy ra $\sin \left(\alpha  +\beta \right)=\sin \alpha\cdot\cos \beta +\cos \alpha\cdot\sin \beta =\dfrac{1}{3}\cdot\left(-\dfrac{2}{3}\right)+\left(-\dfrac{2\sqrt{2}}{3}\right)\cdot\dfrac{\sqrt{5}}{3}=-\ \dfrac{2+2\sqrt{10}}{9}$.\\
		Vậy $\sin \left(\alpha +\beta \right)=-\ \dfrac{2+2\sqrt{10}}{9}$.} 
\end{ex}


%%%%%%%%%%%%%%%% TRẢ LỜI NGẮN %%%%%%%%%%%%%%%%

% nguồn: Ok_TraLoiNgan_Lop11_Lan1
\begin{ex}%[0D6Y3-1]%[ma:ccdafe66]
	 Cho $\tan x=\dfrac{1}{2}$. Tính $\tan \left(x+\dfrac{\pi}{4}\right)$.
	 \shortans{3}
	\loigiai{Ta có $\tan \left(x+\dfrac{\pi}{4}\right)=\dfrac{\tan x+\tan \dfrac{\pi}{4}}{1-\tan x\cdot \tan  \dfrac{\pi}{4}}=\dfrac{\dfrac{1}{2}+1}{1-\dfrac{1}{2}} = 3$.
	} 
\end{ex}

% nguồn: Ok_TraLoiNgan_Lop11_Lan1
\begin{ex}%[0D6B3-1]%[ma:86974c5c]
	 Biểu thức $ A=\cos20^{\circ} +\cos40^{\circ} + \cos60^{\circ}+\cdots+\cos160^{\circ}+\cos180^{\circ}$ có giá trị bằng\\
	\shortans{-1} 
	 \loigiai{Ta có
	\begin{eqnarray*}
		A&=&(\cos20^{\circ}+\cos160^{\circ})+(\cos40^{\circ}+\cos140^{\circ})+\cdots+(\cos80^{\circ}+\cos100^{\circ})+\cos180^{\circ}\\ 
		 &=&(\cos20^{\circ}-\cos20^{\circ})+(\cos40^{\circ}-\cos40^{\circ})+\cdots+(\cos80^{\circ}-\cos80^{\circ})+\cos180^{\circ} \\
		 &=& 0+ 0 +\cdots+ 0 -1\\
		 &=& -1.
	\end{eqnarray*}
	} 
\end{ex}
